\section{Requisitos}
\subsection{Requisitos Funcionais}

\FR{Pipeline de Ataque}{execucao_da_pipeline_de_ataque}{O sistema deve reproduzir o ataque segundo uma pipeline de instruções.}
\FR{Faseamento da Pipeline}{pipeline_de_ataque_faseada_por_configuracoes}{O sistema deve conter a pipeline de ataques faseada em scenario, workload dataset, scenario-specific metric, algorithm (LLM), linguagem de programação e attack load.}
\FR{Selecionar Configuração das Fases}{selecionar_configuracao_fases}{O sistema deve permitir selecionar a configuração de cada fase da pipeline através de configurações pré-definidas.}
\FR{Inserção de Configurações Próprias}{configuracao_propria}{O sistema poderia opcionalmente aceitar inserção por parte do utilizador de configurações de cada fase.} %TODO: Deve ser possível adicionar novas LLMs sem modificar alterar a framework. / Deve ser possível adicionar novos tipos de ataque ou fault loads sem alterações estruturais./ uploaded de ataques diretamente por utilizadores


% \FR{Seleção de LLMs}{selecao_llms}{O sistema deve permitir a seleção de um ou mais modelos de LLMs para serem testados/avaliados.}
\FR{Histórico de Resultados}{historico_resultados}{O sistema deve permitir visualizar resultados de ataques anteriormente executados pelo utilizador.}
\FR{Reutilização de Resultados}{reutilizacao_resultados}{O sistema deve questionar o utilizador quanto à reutilização dos resultados de ataques anteriormente executados no sistema, casos estes estejam públicos.}
\FR{Filtragem de Resultados}{filtragem_resultados}{O sistema deve permitir filtrar os resultados de ataques por scenario, workload dataset, scenario-specific metric, algorithm (LLM), linguagem de programação e attack load.}
\FR{Comparar Resultados}{comparar_resultados}{O sistema deve permitir comparar lado a lado resultados de diferentes ataques, até 3 simultâneo.}

\FR{Registo de Utilizador}{registo_de_utilizador}{O sistema deve permitir a criação de contas para cada utilizador.}
\FR{Login de Utilizador}{login_de_utilizador}{O sistema deve permitir que o utilizador efetue o login.}
\FR{Controlo de Acesso aos Resultados}{controlo_acesso_resultados}{O sistema deve permitir ao utilizador disponibilizar os resultados dos experimentos de forma pública ou privada.}

\FR{Análise do Código de Linguagens de Programação Comuns}{analise_do_codigo_de_linguagens_de_programação_comuns}{O sistema deve permitir compilar e avaliar através de Static Analysis Tools  as vulnerabilidades do código de linguagens de programação comuns (e.g. python, java, javascript, etc).}
\FR{Suporte de Diferentes Tipos de Dados}{suporte_diferentes_tipos_de_dados}{O sistema poderia opcionalmente efetuar ataques com a integração de outros tipos de dados para além de texto (e.g. mensagens em imagens, audios).}
\FR{Suporte a Ataques Diferentes}{suporte_ataques_diferentes}{O sistema poderia opcionalmente efetuar ataques que utilizem code-to-code e code completion prompts, além de ataques de ataques text-to-code.}
\FR{LLM Atacante}{llm_atacante}{O sistema deve implementar metedologias de ataques baseados em LLM de forma a criar prompts adversariais dinamicamente.}

\FR{Armazenamento de Dados}{armazenamento_dados}{O sistema deve permitir armazenar os dados de entrada, a resposta dos modelos e as respetivas métricas.}
\FR{Cálculo de métrica ASR}{calculo_metrica_asr}{O sistema deve permitir calcular o Attack Success Rate (ASR), medindo a proporção de prompts adversariais que burlam com sucesso os mecanismos de segurança.}
\FR{Cálculo de métrica ORR}{calculo_metrica_orr}{O sistema deve permitir calcular o Over-Refusal Rate (ORR), medindo a proporção de prompts benignas que são incorretamente rejeitadas.}
\FR{Cálculo de métrica AOR}{calculo_metrica_aor}{O sistema deve permitir calcular o Achieved Objective Rate (AOR), medindo a amplitude do sucesso do ataque através da proporção de objetivos distintos que foram alcançados.}
\FR{Mecanismo de Júri}{mecanismo_juri}{O sistema deve ter um mecanismo de julgamento através de LLMs para determinar se a resposta do modelo alvo constitui uma recusa em comprir o objetivo, utilizando sistema de votação entre os elementos do júri (LLMs).}

\FR{Avaliação da Utilidade}{avaliacao_utilidade}{O sistema poderia opcionalmente através do mecanismo de júris determinar se a resposta do modelo alvo constitui uma resposta útil.}

\FR{Captcha}{captcha}{O sistema poderia opcionalmente a cada login e registo na plataforma exigir verificação Captcha para verificar se não é se o user não é um Bot.}


% \FR{C}

% | Tipo de requisito               | Redação recomendada                   | MoSCoW      | Semântica                                              |
% | ------------------------------- | ------------------------------------- | ----------- | ------------------------------------------------------ |
% | **Essencial**                   | O sistema **deve permitir** [...]     | Must have   | Obrigatório; sem isto o sistema não cumpre o propósito |
% | **Importante, mas não crítico** | O sistema **deverá permitir** [...]   | Should have | Desejável; adiado se houver restrições                 |
% | **Opcional**                    | O sistema **pode permitir** [...]     | Could have  | Facultativo; apenas se sobrar tempo/recursos           |
% | **Excluído**                    | O sistema **não deve permitir** [...] | Won’t have  | Fora do escopo ou adiado para outra fase               |



\subsection{Requisitos Não Funcionais}

\NFR{Segurança}{O sistema deve implementar autenticação segura com hashing de passwords e e-mails (bcrypt ou superior) para proteger contas de utilizadores.}


\NFR{Usabilidade}{O sistema deve permitir que utilizadores configurem e iniciem um ataque básico em menos de 5 minutos, utilizando configurações pré-definidas.}

\NFR{Compatibilidade}{O sistema deve funcionar nos principais navegadores web (Chrome, Firefox, Safari) nas suas versões mais recentes.}

\NFR{Conformidade}{O sistema deve permitir aos utilizadores exportar e eliminar os seus dados pessoais, cumprindo requisitos básicos do GDPR.}

%TODO: maybe implementar Captcha para dar log in e registo


\NFR{Usabilidade}{O sistema deve permitir que utilizadores configurem e iniciem um ataque básico em menos de 5 minutos, utilizando configurações pré-definidas.}

\NFR{Observabilidade}{O sistema deve registar operações críticas (login, execução de ataques, erros) com timestamp e ID de utilizador para auditoria.}
